%% ma: L12-B06-0076
%% lop: 12
%% chuong: 2
%% bai: 6
%% dang: 12.6.13
%% loai: TN4
%% muc: TH
%% dap_an: "A"
%% nguon: "CĐ2.1.pdf giáo viên gửi 27/10/2026 (M07, Lê Minh Khôi – Nguyễn Anh Khôi, Chương II Chủ đề 1 Vectơ trong không gian), Câu 37"
%% kien_thuc: "Bài 6 (bình phương vô hướng)"
%% trang_thai: chinh-thuc
\begin{ex}
Cho hai vectơ $\vec a$ và $\vec b$. Đặt $A=\vec a\cdot\vec b$ và $B=\left|\vec a-\vec b\right|^2-\left|\vec a+\vec b\right|^2$. Biết rằng với mọi $\vec a$ và $\vec b$ thì $A=kB$ với $k\in\mathbb R$. Giá trị của $k$ bằng
\choice{\True $-0{,}25$}{$0{,}25$}{$-0{,}50$}{$0{,}50$}
\loigiai{$B=\left(\vec a^2-2\vec a\cdot\vec b+\vec b^2\right)-\left(\vec a^2+2\vec a\cdot\vec b+\vec b^2\right)=-4A$, nên $k=-0{,}25$.}
\end{ex}
